\clearpage
% ============================================================================
\section*{v4.10.0-r72 Source-Only: The Riemann Non-Claim as a Dormant Note; the Photon-Clock Guard without RH}
\addcontentsline{toc}{section}{v4.10.0-r72 Source-Only: The Riemann Non-Claim as a Dormant Note; the Photon-Clock Guard without RH}
\label{sec:v4100-r72-rh-dormant}
% ============================================================================

\paragraph{Purpose.}
Since v4.8.0-r04 (Section~\ref{sec:v480-r04-riemann-decoupling}) no result of the Companion rests on the Riemann source:
the compensator, the cascade and the triole hierarchy have standalone definitions, and every other dependency was reclassified.
Nevertheless the guard \(W_\gamma^{\rm glob}\) has been listed in every later layer as ``a global photon-clock theorem''.
The Riemann hypothesis was never a goal of the Companion, and repeating the non-claim in every layer gives it a weight it does not have.
This layer records the non-claim once, as a dormant note, and removes the Riemann hypothesis from the list of active guards.
\textbf{Editorial layer}: no statement, test, number or status of an open problem changes; there is no H~number and no computation.

\noindent\textbf{Source-only status.}
This v4.10.0-r72 layer is source-only.  It does not build a PDF, does not assign a new DOI,
and does not authorize a public release.  The v5.5.0 Version DOI
\href{https://doi.org/10.5281/zenodo.24563811}{10.5281/zenodo.24563811} and the Concept DOI
\href{https://doi.org/10.5281/zenodo.19210728}{10.5281/zenodo.19210728} are unchanged by it.

\subsection*{v4.10.0-r72 the dormant non-claim}
\label{subsec:v4100-r72-non-claim}

\paragraph{Remark v4.10.0-r72.1 (dormant Riemann non-claim).}
The Companion makes no claim about the Riemann hypothesis or about the zeros of the zeta function.  The Riemann source
\cite{brendecke2026rh} is not a proof source (v4.8.0-r04); the rows that drew on it are withdrawn or archived there.
This note is \emph{dormant}: it is not an open premise, it is spent by no layer, and it need not be repeated.  It becomes active
again only if a later layer reintroduces a link to the Riemann zeros; such a layer must then state the link and its status explicitly.

\paragraph{Remark v4.10.0-r72.2 (the photon-clock guard without RH).}
The fourth component of the guard vector keeps its content and its value.  From this layer on it is named
\(W_\gamma^{\rm glob}\): \emph{no global photon-clock theorem}.  The photon-clock construction remains a local handoff probe and
selector window (v4.7.0 development chapter); what the guard forbids is its promotion to a global theorem.  The Riemann
hypothesis is no longer part of its name, because no layer since v4.8.0-r04 uses the Riemann zeros as a clock.

\paragraph{Remark v4.10.0-r72.3 (closing formula).}
From this layer on, the closing sentence of a layer reads: ``The layer does not imply exact finite \(C5\) zero closure, exact
all-readout terminality, \(C6/T2\), a global photon-clock theorem, or global all-readout closure.''  The guard vector
\(\Gamma_{\rm guard}=(1,1,1,1,1)\) is unchanged; only the name of its fourth component is shorter.

\paragraph{Remark v4.10.0-r72.4 (editorial note on earlier layers).}
By the author's decision of 27 November 2026 the same closing formula in the earlier layers of this source (v4.7.0 development
chapter, v4.8.0, v4.9.0 and v4.10.0 layers, 193 places) now reads ``a global photon-clock theorem'' instead of ``a global
photon-clock/RH theorem''.  This is a deliberate and recorded exception to the append-only rule, limited to these two words of a
recurring formula: no statement, number, verdict or guard value of an earlier layer changes.  The places where earlier layers
explain the withdrawal of the Riemann source (v4.8.0-r04 and the status notes that cite it) are kept unchanged.  The published
versions on Zenodo keep their original wording.

\subsection*{v4.10.0-r72 frames}
\label{subsec:v4100-r72-frames}

No frame quantity occurs.  Convention~F is not touched.

\subsection*{v4.10.0-r72 register entries}
\label{subsec:v4100-r72-register}

\begin{center}
\small
\begin{tabular}{@{}>{\raggedright\arraybackslash}p{0.40\textwidth}>{\raggedright\arraybackslash}p{0.52\textwidth}@{}}
\toprule
statement & status after v4.10.0-r72\\
\midrule
Riemann non-claim (r72.1) & dormant note; not an active guard; Riemann source no proof source since v4.8.0-r04\\
guard \(W_\gamma^{\rm glob}\) (r72.2) & active, value unchanged; name ``no global photon-clock theorem''\\
closing formula of later layers (r72.3) & without RH\\
OP1--OP5, C5, C6/T2 & open, as before\\
\bottomrule
\end{tabular}
\end{center}

\subsection*{v4.10.0-r72 notes}
\label{subsec:v4100-r72-notes}

{\raggedright
Notes: \texttt{notes/\allowbreak research/\allowbreak v4\_\allowbreak 10\_\allowbreak 0\_\allowbreak r72\_\allowbreak riemann\_\allowbreak ruhend/}.
Tools: none (editorial layer, no computation).  No Wolfram, no Qiskit, no HDL/MMA comparison.
\par}

\subsection*{v4.10.0-r72 ledger and guard update}
\label{subsec:v4100-r72-ledger-guard}

\[
R_{\rm chk}^{v4100\text{-}r72}=(1,0,0)
\]
for (status recorded; no computation; no comparison with data).

The layer spends no guard:
\[
R_{\rm glob}^{(v4100\text{-}r72,+)}=(0,1,1,1,1),
\qquad
\Gamma_{\rm guard}^{v4100\text{-}r72}=(1,1,1,1,1),
\qquad
\beta_4=1/256\ \text{unspent}.
\]
The layer does not imply exact finite \(C5\) zero closure, exact all-readout terminality, \(C6/T2\), a global photon-clock theorem,
or global all-readout closure.  OP1--OP5, C5 and C6/T2 remain open.
